\subsection{Probleme und Anpassungen}
Aufgrund der neuen GitOps-Verwaltung musste das Autoscript zum Erstellen des Kind-Clusters angepasst werden.
Nun wird, anstatt dass alles automatisch in den Cluster geladen wird, zuerst FluxCD auf den Kind-Cluster installiert.
Darauf folgt die gleiche Reihenfolge wie beim Prod-Cluster. Die Git-Repository-Konfiguration und das Root-Kustomization für die Developer-Umgebung befinden sich im Ordner „./.k8s” und werden jeweils reingeladen.
Die Zugangsdaten, die FluxCD benötigt, um auf das GitOps-Repository zugreifen zu dürfen, müssen in einer .env-Datei eingetragen werden. Dabei nutzt jeder Studierende seinen eigenen Username und seinen eigenen Access-Token.
Anschließend richtet FluxCD auch den lokalen Cluster ein.

Ein Problem entstand dadurch, dass Codetask in der öffentlichen Instanz von GitLab liegt und dort in der kostenlosen Variante keine Projekt-Access-Token verwendet werden können.
Daher konnten diese nicht als Zugangsdaten von FluxCD für das GitOps-Repository verwendet werden.
Auch die Verwendung eines SSH-Keys war leider nicht möglich, da der SSH-Port 22 blockiert wird, wenn zu GitLab.com ein Known Host erstellt werden soll.
Auch ein Deploy-Token konnte nicht genutzt werden, da FluxCD nicht nur Lese-, sondern auch Schreibrechte auf das Repository benötigt.
Daher wurde als Übergangslösung ein „Key Bot” erstellt. Dabei handelt es sich um einen angelegten User mit Maintainer-Rechten in der Gruppe Kompetenzquartet, dessen Person Access Token benutzt wird.

Ein großes Problem ergab sich dadurch, dass Codetask nicht mit mehreren Replicas der Datenbank umgehen konnte, wodurch sich die Anwendung asynchron verhielt.
Um dies zu verbessern, musste zunächst der Connection String zur Datenbank am Ende um „\&replicaSet=rs0” ergänzt werden, damit MongoDB Replicas verwalten kann.
Zum anderen wurde das Helm-Chart, in dem das StatefulSet der MongoDB war, durch einen Percona-Operator ersetzt. Operatoren werden in Kubernetes verwendet, um wiederkehrende Aufgaben, die meist manuell erfüllt werden müssen, zu automatisieren.
Dabei überwacht er den gewünschten Zustand einer Anwendung und führt anwendungsspezifische Logik in Kubernetes aus.
In diesem Fall übernimmt Percona die Einrichtung und Verwaltung von MongoDB mit einem Replica-Set.
Zum Vergleich: Ein StatefulSet ist nur dafür verantwortlich, Pods zu verwalten. Ein Percona Operator verwaltet hingegen nicht nur die Pods, sondern zusätzlich auch fachlogisch Mongodb.
